\section{Bounded-factor nesting estimates}\label{sec:uniform-nesting}

The exponent-level saturation proved above identifies a nonzero normalization. We now strengthen both the comparison between plane and cylinder and the analytic estimate. The first removes large polygons at bounded cost; the second proves convergence of the normalized edge contributions. Together they prove the first assertion of Theorem~\ref{thm:nesting}.

\subsection{Uniform geometric comparison}\label{c6:uniform-geometry}

The height tails of the preceding section discard polygons much larger than the period, but their polynomial prefactor does not yet give a bounded cost at scale $N$. We first bound the individual polygon mass in each planar diameter annulus. A refinement of the cylinder free-energy estimate will then leave only a uniformly bounded number of exceptional annuli.

Keep $\beta=3/8$. Let $D$ be a finite cylinder between flat caps, or a finite convex tiled polygon in the plane. Write $Q_D=Z_H(\sqrt2)$ in the cylinder and $Q_D=1$ in the plane, and use the same notation for restrictions of the available essential polygons. All these polynomials are positive. For a set $K$ of occupied centers put
\[
 H_D(K)=Q_{D\setminus K}/Q_D,
\]
where the numerator removes the essential polygons hitting $K$. These ratios increase as the available polygon set grows. Indeed, for one deleted polygon $v$ in a set $P$,
\[
 \frac{Q_P}{Q_{P\setminus v}}
 =1-t\kappa^{|v|}\frac{Q_{P\setminus N[v]}}{Q_{P\setminus v}},
\]
where $N[v]$ is its closed conflict neighborhood. Induction on $|P|$, deleting its neighbors successively in the smaller set, gives the claim; extra neighbors only increase the deletion ratio. Products give the assertion for several deletions. The same reasoning applies for $0\le t\le\sqrt2$.

Write $\mathcal Y(y)$ for the partition with these essential polygons and, in addition, mutually disjoint contractible polygons enclosing the triangular lattice vertex $y$, each with factor $2$. Expanding in the enclosing polygons leaves positive restricted essential partitions, so $\mathcal Y(y)\ge Q_D$, including when $y$ lies on a cap.

For a simple polygon $L\subset D$ and a boundary source $a$, let $P(a,L)$ be the real part of the normalized signed mass of stems first hitting $L$. Stop a stem at the incoming midpoint before visiting its first center of $L$; include all essential polygons disjoint from the stem and $L$, and divide by $Q_{D\setminus L}$. The contour identity identifies this with $1$ minus the real normalized mass exiting to the original boundary in the domain with $L$ removed. This convention gives $P(a,L)=1$ if the first center towards which $a$ points has been removed. Since outer-exit real phases are positive and the deletion ratios increase before removal,
\begin{equation}\label{eq:first-hit-comparison}
 1\ge P(a,L)\ge
 \sum_{\substack{\gamma:\,a\to\partial D\\\gamma\cap L\ne\varnothing}}
 \kappa^{|\gamma|}\cos(\beta W(\gamma))H_D(\gamma)\ge0.
\end{equation}

\begin{auxlemma}[A primitive of the boundary flux]
Choose a dual-edge chain $\Gamma$ from a boundary lattice vertex $y_0$ to an interior lattice vertex $y$. For an oriented honeycomb link, count a crossing of $\Gamma$ positively when the chain goes from its right face to its left face. Let $\operatorname{cr}_\Gamma(\gamma)$ be the sum of these signs along a completed path and set
\[
 I=\Im\sum_{a\in\partial D}\sum_{\gamma:\,a\to\partial D}
   \operatorname{cr}_\Gamma(\gamma)\,
   \kappa^{|\gamma|}e^{i\beta W(\gamma)}H_D(\gamma).
\]
The sum over $a$ is over all ordinary boundary ports. Then
\begin{equation}
 2\!\!\sum_{L\text{ contractible enclosing }y}
       \kappa^{|L|}H_D(L)\sum_aP(a,L)
 \le |I|+\sqrt2\!\sum_{L\text{ essential}}
       \kappa^{|L|}H_D(L)\sum_aP(a,L).
 \label{eq:legacy-P-G1}
\end{equation}
The essential sum is absent in the plane.
\end{auxlemma}
\begin{proof}
We compute the same prefix current in two ways. For a prefix $\eta$ from $a$ to an oriented internal-link midpoint $e$, write $\epsilon_\Gamma(e)$ for the crossing sign just specified, or zero if the chain does not cross $e$. Define
\[
 J=\Im\sum_{a\in\partial D}\sum_{e,\eta:\,a\to e}
     \epsilon_\Gamma(e)\,
     \kappa^{|\eta|}e^{i\beta W(\eta)}H_D(\eta).
\]
Only centers already visited by $\eta$ are deleted in $H_D(\eta)$. Repeated chain crossings, if present, are counted with their signed multiplicities. Figure~\ref{fig:geometry-current} records the orientation convention.
\begin{figure}[htbp]
\centering
\begin{tikzpicture}[x=.95cm,y=.95cm,>=Latex,font=\small,
  primal/.style={draw=blue!65!black,very thick},
  dual/.style={draw=orange!80!black,thick,densely dashed}]
  \path[draw=black!40,fill=blue!6]
    (0,-.69282)--(0,.69282)--(-1.2,1.38564)--(-2.4,.69282)
    --(-2.4,-.69282)--(-1.2,-1.38564)--cycle;
  \path[draw=black!40,fill=orange!7]
    (0,-.69282)--(1.2,-1.38564)--(2.4,-.69282)--(2.4,.69282)
    --(1.2,1.38564)--(0,.69282)--cycle;
  \draw[primal,->] (0,-.69282)--(0,.69282);
  \node[anchor=south] at (0,.90) {$e$};
  \draw[dual,->] (1.2,0)--(-1.2,0);
  \fill (-1.2,0) circle (1.8pt) node[below=4pt] {$y$};
  \fill (1.2,0) circle (1.8pt) node[below=4pt] {$z$};
  \node at (-1.2,.68) {left face};
  \node at (1.2,.68) {right face};
  \node at (0,-1.67) {crossing sign $+1$};
  \node at (0,-2.3) {(a) Dual edge $[y,z]$, traversed from $z$ to $y$};

  \begin{scope}[xshift=6.65cm]
    \path[fill=blue!6]
      (1.5,-.7)--(1.5,.7)--(0,1.55)--(-1.5,.7)
      --(-1.5,-.7)--(0,-1.55)--cycle;
    \draw[primal]
      (1.5,-.7)--(1.5,.7)--(0,1.55)--(-1.5,.7)
      --(-1.5,-.7)--(0,-1.55)--cycle;
    \draw[primal,->] (1.5,-.45)--(1.5,.5);
    \draw[primal,->] (-.33,1.363)--(-.95,1.012);
    \node[anchor=south] at (0,1.82) {$L$ (counterclockwise)};
    \draw[black!45] (2.7,-1.65)--(2.7,1.65)
      node[above] {$\partial D$};
    \draw[dual,->] (2.7,0)--(0,0);
    \fill (2.7,0) circle (1.8pt) node[right=3pt] {$y_0$};
    \fill (0,0) circle (1.8pt) node[below=4pt] {$y$};
    \node at (0,-.78) {$\operatorname{cr}_\Gamma(L)=1$};
    \node at (.4,-2.3) {(b) Net crossing equals enclosure};
  \end{scope}
\end{tikzpicture}
\caption{The crossing convention for the current primitive.
(a) The oriented honeycomb link $e$ crosses the tiling edge $[y,z]$;
the dual traversal from its right face to its left face has sign $+1$.
Reversing either orientation reverses this sign.
(b) For a counterclockwise contractible polygon, the interior lies to
the left. Hence a dual chain from the exterior boundary gap $y_0$ to
$y$ has net crossing $1$ if the polygon encloses $y$, and $0$ otherwise.
The global polygon and chain are schematic. These are topological
crossing signs; the imaginary prefix-current coefficient is computed
separately in the proof.}
\label{fig:geometry-current}
\end{figure}


\paragraph{Reverse at the cut link.}
Reverse the prefixes ending at $e$ and apply the cut-link calculation from Proposition~\ref{prop:occupancy}. Permit contractible detached polygons about the cut midpoint with factor $2$, and retain the essential detached polygons. The essential blocked-loop turns are still $\pm2\pi/3$, irrespective of the stem side, so their cancellation is unchanged. A return through the cut along an essential polygon has turn zero. A contractible return has turn $\pm2\pi$; reversal conjugates the phase and exchanges the two link directions. Thus the imaginary current across a dual edge is $-1/\sqrt2$ times the left-minus-right difference of the centered partitions, divided by $Q_D$. Each enclosing polygon crossing that dual edge is counted once with its factor $2$. Along outer exits from the cut no enclosing contractible polygon remains. Telescoping from $y_0$ to $y$ therefore gives
\begin{equation}\label{eq:prefix-current-potential}
 J=-\frac{\mathcal Y(y)/Q_D-1}{\sqrt2}.
\end{equation}
Here $\mathcal Y(y_0)=Q_D$ because a contractible polygon in $D$ cannot enclose a boundary vertex.

\paragraph{Continue the prefixes.}
Instead continue each prefix to its first exit, treating its past centers as additional deleted boundary. The starting residual partition is exactly the one already attached to the prefix. A loop in this continuation cannot surround its source, which remains accessible along the past. The exits are either completed boundary paths, giving $I$, or blocked returns to a past center. Each blocked return consists uniquely of a first-hitting stem followed by an incomplete simple polygon; all detached polygons avoid both pieces.

Pair the two orientations of this polygon, including their opposite link-crossing signs. The coefficients of their circulation are
\[
\begin{array}{c|c|c}
 \text{polygon}&\text{additional turns}&\text{oriented phase difference}\\ \hline
 \text{contractible}&\pm4\pi/3&e^{i\pi/2}-e^{-i\pi/2}=2i\\
 \text{essential}&\pm2\pi/3&\pm(e^{i\pi/4}-e^{-i\pi/4})=\pm i\sqrt2.
\end{array}
\]
For a contractible polygon, the counterclockwise integrated crossing equals its enclosure indicator at $y$. The imaginary part of $2i$ times its signed stem weight is twice the real stem weight, namely the factor $P(a,L)$ after normalization. An essential polygon has integrated crossing of modulus at most one. Consequently, with
\[
 A=\sum_{L\text{ contractible enclosing }y}
       \kappa^{|L|}H_D(L)\sum_aP(a,L),
\]
this second computation reads
\begin{equation}\label{eq:prefix-current-circulation}
 J=I+2A+E,\qquad
 |E|\le\sqrt2\sum_{L\text{ essential}}
       \kappa^{|L|}H_D(L)\sum_aP(a,L).
\end{equation}
Combining the two computations gives
\[
 2A+\frac{\mathcal Y(y)/Q_D-1}{\sqrt2}=-I-E.
\]
The additional term on the left is nonnegative, proving \eqref{eq:legacy-P-G1}.
\end{proof}

\paragraph{A planar annular bound.}
For $s\ge1$, take a regular lattice hexagon of inradius comparable to $s$ about $y$. Only chords separating $y$ from $y_0$ contribute to $I$, with net crossing of modulus one. From a source $a$ such a chord has diameter at least a constant times $\min(s,|a-y_0|)$. Summing the boundary diameter estimate over the ports gives $|I|\le Cs^{3/4}$.

Enlarge the hexagon by a fixed factor. A polygon enclosing $y$ with diameter in $[s,2s]$ lies inside it and has transverse span at least $cs$ in one pure direction. For order $s$ sources on the corresponding cap, choose parallel tubes of width a sufficiently small multiple of $s$, with axes in the central part of that transverse span. The polygon has a subarc crossing each tube from one lateral side to the other. Every confined path between the two caps in that tube must therefore hit the polygon. Lemma~\ref{lem:narrow-channel-bridges} and \eqref{eq:first-hit-comparison} give $\sum_aP(a,L)\ge cs^{3/4}$. Applying \eqref{eq:legacy-P-G1} in the plane yields
\begin{equation}
 \sum_{L\supset_{\rm int}y:\,\operatorname{diam}(L)\in[s,2s]}
       \kappa^{|L|}\le C.
 \label{eq:legacy-P-G2}
\end{equation}

\paragraph{Cap-source bounds on a cylinder.}
If $y$ is at distance at least a fixed positive multiple of $N$ from both caps, the same current satisfies $|I|\le CN^{3/4}$. Bridges have real phase and give no imaginary contribution. A same-cap path separates the cylinder after closure along its cap, so its net contribution has modulus at most its mass and requires lifted diameter at least a constant times $\min(N,\operatorname{dist}(a,y_0))$.

For $r\le c_1N$, the real normalized completed mass from one cap source with diameter above $r$ is $O(r^{-1/4})$. Indeed its complement contains the base returns in a centered triangle of size comparable to $r$, with deletion ratios at least one; subtract their mass from the positive contour identity. The same comparison in the domain with $L$ removed gives
\begin{equation}\label{eq:cap-first-hit}
 P(a,L)\le C\bigl(1+\min(N,\operatorname{dist}(L,\text{cap of }a))\bigr)^{-1/4}.
\end{equation}
These bounds will control the source sums in \eqref{eq:legacy-P-G1}.

\begin{auxlemma}[All but boundedly many dyadic heights]
For each $u>0$ there is $C_u<\infty$, independent of $N$, such that all but at most $C_u$ heights $H_j=d_0N2^j$, $j\ge0$, satisfy
\begin{equation}
 \sum_{L\text{ essential in the }H_j\text{-slab}}
       \kappa^{|L|}H_{\rm slab}(L)\le e^{uH_j/N}.
 \label{eq:legacy-P-g}
\end{equation}
The property is unchanged by admissible translation of the caps. We call these heights good; the exceptional indices need not lie in a bounded interval.
\end{auxlemma}
\begin{proof}
Put $F_H(t)=-\log Z_H(t)$. Restricting to disjoint subslabs and using deletion monotonicity shows that $F_H$ is superadditive in height. The free-energy bound \eqref{eq:cylinder-free-energy} makes
\[
 a_j=\frac{N F_{H_j}(\sqrt2)}{H_j}
\]
a nondecreasing sequence in $[0,C]$, apart from a fixed initial range which can be included among the exceptions. Fix a sufficiently large dyadic integer $k=2^r$, depending on $u$. Since
\[
 \sum_{j\ge r}(a_j-a_{j-r})\le rC,
\]
at most $4rC/u$ indices violate
\[
 F_H(\sqrt2)-kF_{H/k}(\sqrt2)\le uH/(4N).
\]
This proves the required uniform count at the first step.

For $h=H/k\gg N$, \eqref{eq:cylinder-integrated-ratio} bounds the unadorned essential-polygon mass by $Ch/N$. Restricted partition polynomials are at most $1$, and $Z_h(t)\ge e^{-Ch/N}$. Hence
\[
 0\le F'_h(t)\le C(h/N)e^{Ch/N},\qquad 0\le t\le\sqrt2.
\]
For $t\in[\sqrt2-e^{-2Ch/N},\sqrt2]$, superadditivity and the preceding increment bound give
\[
 F_H(\sqrt2)-F_H(t)\le uH/(3N)
\]
after adding a bounded initial set of indices to the exceptions. Each deletion-ratio numerator increases as $t$ decreases. Every summand of $F'_H(t)$ in this interval is therefore at least its endpoint value times $e^{-uH/(3N)}$. Integrating and using \eqref{eq:cylinder-integrated-ratio} gives
\[
 \sum_{L\text{ essential}}\kappa^{|L|}H_{\rm slab}(L)
 \le C(H/N)e^{(u/3+2C/k)H/N}\le e^{uH/N}.
\]
Choose $k$ large enough and again include the bounded initial range among the exceptions. All bounds are invariant under row translations.
\end{proof}

\paragraph{Contractible polygons at a good height.}
Choose $u>0$ small enough for the barrier comparison below. For a good height $H$, contractible polygons enclosing $y$ with vertical span in $[H/4,H/2]$ have total mass at most $Ce^{-cH/N}$. To prove this, bin their minimum levels into intervals of width $\delta N$. Since each polygon encloses $y$, there are $O_\delta(1+H/N)$ bins. For one bin take an encompassing slab of height $H$, starting just below it. Inside this slab vary the two caps of a smaller containing cylinder in windows of length comparable to $\delta N$, so each polygon remains at a lower-cap distance between $\delta N$ and $5\delta N$.

Apply \eqref{eq:legacy-P-G1} and average over the caps. On its right, enlarge the deletion ratios to those of the encompassing slab. For each included essential polygon, the averaged sum of $P(a,L)$ over cap sources is $O_\delta(N^{3/4})$ by \eqref{eq:cap-first-hit}; give excluded polygons contribution zero. The good-height bound controls the remaining essential sum by $e^{uH/N}$, while $|I|=O_\delta(N^{3/4})$.

On the left, the slab barriers give each polygon a deletion ratio at least $ce^{c_0H/N}$. Its source sum is at least $c_\delta N^{3/4}$. For the latter assertion use the confined upper arches from the slab construction, between ports near horizontal offsets $-N/4,+N/4$ from a minimum point of the polygon. Their separations lie in $N/2\pm\epsilon N$. For sufficiently small $\delta$, the part of each arch below height $6\delta N$ stays within $N/10$ horizontally of these two endpoints: the initial triangle is small, the tilted tube rises above that height after its deep exit, and all subsequent pieces stay higher. Closing along the base interval therefore surrounds the minimum point but not the maximum, forcing intersection. The construction gives mass $\gtrsim N^{-5/4}$ on order $N^2$ endpoint pairs, hence the stated source sum by \eqref{eq:first-hit-comparison}.

Taking $u<c_0$ bounds the polygon mass in the bin exponentially. The number of bins is polynomial in $H/N$ and is absorbed by reducing the exponential rate. Thus
\begin{equation}\label{eq:good-contractible-span}
 \sum_{\substack{L\text{ contractible enclosing }y\\
        \operatorname{span}_{\rm vert}(L)\in[H/4,H/2]}}
       \kappa^{|L|}\le Ce^{-cH/N}
\end{equation}
for every good sufficiently large height.

\begin{auxproposition}[Bounded cost of large specified lifts]
Fix $c_2>0$. On a cylinder of period $N$, the total individual critical weight of contractible polygons enclosing a fixed mark $y$ whose specified lifts about a fixed preimage of $y$ have diameter greater than $c_2N$ is bounded by $C_{c_2}$, independently of $N$.
\end{auxproposition}
\begin{proof}
The good vertical-span bands have summable total weight by \eqref{eq:good-contractible-span}. A uniformly bounded number of exceptional bands remain, together with the band of vertical spans $O(N)$.

Consider one exceptional band with upper scale $S\gtrsim N$. Its original vertical span is both at most and at least fixed positive multiples of $S$. For the initial $O(N)$ band put $S\asymp N$ instead; its lift diameter is bounded below by $c_2N$. Choose a lattice period tilted by $60^\circ$, of length comparable to $S$ and with vertical component greater than twice the original span upper bound. The specified lift and its enclosed disk lie in a horizontal strip of that bounded width. They therefore project injectively onto this tilted cylinder, since distinct period translates of the strip are disjoint. The image is contractible and determines the original lift uniquely as the lift whose disk contains the prescribed preimage of $y$. This map preserves weights and is injective on the class being counted.

Write $s'$ for the span normal to the tilted period and $d$ for the lift diameter. If $d$ is a sufficiently large multiple of $S$, then $s'\asymp d$: two nonparallel projections control diameter, and the original vertical projection is $O(S)$. Apply the good-height estimate on the tilted cylinder. Its good $s'$-bands again have summable weight. It has only a uniformly bounded number of exceptional $s'$-bands, each corresponding to a diameter interval with bounded ratio of endpoints, so \eqref{eq:legacy-P-G2} bounds each by a constant.

The remaining class $s'=O(S)$ has $d=O(S)$. In an original exceptional band it also has $d\gtrsim S$ by the original lower span bound. In the initial band its lower bound $d>c_2N$ is likewise a fixed multiple of $S$. Thus this class is covered by a bounded number of planar diameter annuli, with constants depending only on $c_2$. Summing over the uniformly bounded number of original exceptional bands proves the proposition. For the finitely many periods below the confinement threshold, a simple polygon in a slab of height $s$ uses at most $O(N(1+s))$ centers, so its lift diameter is bounded by the same quantity. The cylinder height tails and the finite bounded-height graphs therefore cover these periods after increasing the constant.
\end{proof}

Discarding polygons from a positive separating-polygon partition costs at most the exponential of their total individual weight times the loop factor: simply forget all avoidance restrictions involving a discarded polygon. Every essential separator meets the horizontal segment between the antipodal marks, so the essential-polygon bound applies to this family. Together with the proposition, it therefore allows us to discard all essential polygons and all contractible lifts of diameter greater than $c_2N$ at bounded multiplicative cost, for any fixed bounded positive loop factor. The planar annular bound also makes a fixed change of $c_2$ harmless. The remaining polygons lift to local systems about one of the two antipodal marks. For loop factor $2$, this is the plane/cylinder comparison written explicitly in \eqref{eq:disk-cylinder-comparison} below, where we specify the disk convention.

\subsection{Uniform analytic comparison}\label{c6:uniform-analytic}

Return to the notation of Section~\ref{c6:gas}; in particular $I=\int_{\mathbb R}V(s)\rho(s)\,ds$ again denotes the bulk constant, not the complete-path current of the preceding subsection. We now sharpen the analytic estimate at the two values needed for the physical ratio. The precise output is convergence of normalized center values along finitely many residue classes. Saturation will make each of their limits positive.

\begin{auxproposition}[Classwise convergence of the normalization]
For $r\in\{0,i\pi/2\}$ put
\[
 \sigma(r)=\frac38+\frac{2r^2}{3\pi^2},\qquad
 A_m(r)=e^{-Im^2}m^{-\sigma(r)}P_m(e^r).
\]
There are a fixed integer $Q\ge1$, a number $\delta>0$, and positive constants $\mathfrak a_b(r)$, $0\le b<Q$, such that, as $m\to\infty$,
\begin{equation}\label{eq:classwise-normalization}
 A_m(r)=\mathfrak a_{m\bmod Q}(r)+O(m^{-\delta}).
\end{equation}
In particular,
\begin{equation}\label{eq:uniform-cylinder-normalization}
 \mathcal Z_{m,m}(1)=P_m(1)/P_m(i)\asymp m^{1/6}.
\end{equation}
By lattice rotation the same bound holds in either other pure period direction.
\end{auxproposition}
\begin{proof}
Use the transformed sum \eqref{eq:gas-grand-sum}, whose exponent is \eqref{eq:legacy-P-star}. Choose a sufficiently large \emph{fixed} $g$ and put
\[
 T=2a_0^{-1}\log m+g+i\theta=T'+i\theta,
 \qquad |\theta|\le\pi/2.
\]
The stability estimate \eqref{eq:legacy-P-P3} and the \hyperref[gas:grand-localization]{localization bound} continue to hold with this choice. The coercivity error is $O(1/g)\mathcal Q$, absorbed by taking $g$ large, and images and contour displacements cost only $O(1)$ per particle in the linear term. Choose the period graph with a fixed cutoff profile in the real distance from either edge $\pm T'/2$. In coordinates relative to $\pm T/2$, its shape is then independent of $m$ on the corresponding half interval.

Multiply the lifted polynomial on this nome circle by
\[
 e^{-2Am^2e^{-2T}}e^{-Im^2}m^{-\sigma(r)},
\]
and denote its value by $F_m(\theta,r)$. The factor $\exp(-2Am^2e^{-2T})$ is analytic in the nome and equals one at its center, so the circle mean is exactly
\begin{equation}\label{eq:normalization-circle-mean}
 A_m(r)=\frac1\pi\int_{-\pi/2}^{\pi/2}F_m(\theta,r)\,d\theta.
\end{equation}
We prove uniform exponentially fast convergence of $F_m$ along each class; the same convergence then holds for this average.

\paragraph{The scalar factor.}
Include $\exp(Tr^2/(4h_*))$ from \eqref{eq:legacy-P-star} with the scalar prefactor \eqref{eq:gas-transform-scalar}. For the quadratic bulk term, retain both the residue at the pole $-a_0i$ and its reflection in the sampling calculation \eqref{eq:gas-bulk-symbol}--\eqref{eq:gas-bulk-period-correction}. After the displayed subtractions, the additional term is a constant multiple of
\[
 m^2e^{-a_0T}=e^{-a_0(g+i\theta)},
\]
with error exponentially small in $T'$. For the logarithmic determinant in \eqref{eq:gas-determinant-rate}, use the complete all-mode sum of $-\log(1-e(p))$ and subtract its value at zero. The symbol is even and regular in a strip, decays exponentially, and continues to sectors of small slope about the real rays. Shifting the inverse Fourier contour and bending its tails therefore makes the nonzero period translates exponentially small. Its contribution is $\pi T/(16L)$ plus a constant and an exponentially small error. The square-root factors have constant product. The powers of $m$ from these terms are exactly $m^{\sigma(r)}$, so the normalized scalar converges to a bounded function of $\theta,r$, uniformly with exponentially small error.

\paragraph{The edge integrals.}
Compare $m$ with any larger $m'$ in the same residue class, at the same $\theta$. In both sums retain only particles whose real positions lie within $\epsilon T'$ of an edge, where $T'$ is the smaller period, retain at most $T'$ particles, and set $j$ equal to minus the total $W$-charge of the right edge. The omitted part has exponentially small integrated absolute weight. A particle away from the edges pays the localization factor $f_0$; large counts are suppressed by the factorials; and any other $j$ leaves $\mathcal D+j\ne0$ throughout a central interval of length comparable to $T'$. Since $g$ is fixed, the bound on the full integrated absolute weight is uniform in $m$.

Parameterize the retained integrals by positions relative to their own period edges. For sufficiently small fixed $\epsilon$, the two integrands agree up to an exponent error $O(e^{-cT'})$. Here are the components of that comparison.

\begin{itemize}

\item 

The \(\alpha WW^t/p^2\) Schur energy combined with \(j\)-terms completes exactly to \(\alpha\int({\cal D}+j)^2 ds\) in the quadratic form, where the staircase follows real order and integration is along the complex graph. Also the \(r\) term involving positions uses only \(\sum W_d s_d+jT\), by the given formula for \(u_*\). Both are unchanged in edge coordinates since \({\cal D}+j=0\) between edges. The residual tilt terms (after extracting the bulk above) involve only counts, \(j\), and kernel constants independent of positions and period.

\item 

In the skew pole keep its sign part and the nearest translate in its smoothing correction, and in the regular symbol keep the nearest translate alone, with the same regularization on self terms. These prescriptions coincide for both sizes (sign of the original difference unchanged); all other images cost exponentially small times squared particle count by the sector decay already proved.

\item 

In each chemical primitive the integral from 0 to the relevant edge \(\pm T/2\) equals \(\pm1/2\) before multiplying by \(i2\pi m\sigma_i\) (even periodized density of mass one, analytically continued). Between there and the particle, \(m\) times the density equals
\[
 (\sqrt3 L)^{-1} e^{-a_0(g+i\theta)/2}
 \left(e^{a_0 v}+e^{-a_0 v}\right)
\]
up to uniform exponentially small additive error, \(v\) the coordinate of integration relative to the edge. This follows by expanding the two nearest images of the sinh quotient, all remaining ones exponentially suppressed; take \(\epsilon\) small and include the fixed offset at the endpoint.

\item

Retain the discrete phase $\varepsilon_m(\nu)$ of each term in \eqref{eq:gas-grand-sum}. It depends on the species counts and winding data, not on the positions or on $T$. The signs were independent of $K_{\rm count}$, so shifting its Poisson sum by the integer $m$ introduces no new sign. The remaining $m$-dependent phases come from fixed logarithm branches raised to the power $m$ and from the half-period chemical primitives. Every branch multiplier is a root of unity: the shifts and offsets are rational multiples of $\pi$, with fixed rational coefficients in the logarithmic factors. The sinh image products, the fixed powers of $i$ from the trigonometric substitutions, and the integer branch jumps give only finitely many types of such coefficients. Write each such multiplier as $e^{2\pi i\rho}$ with $\rho\in\mathbb Q$, and choose $Q$ to be a common multiple of these denominators and of the parity modulus. Multiplying these coefficients by arbitrary integer particle counts or windings introduces no new denominator. Hence the entire discrete phase agrees for $m,m'$ in the same class modulo $Q$. We do not require it to be independent of $m$ across different classes.
\end{itemize}

The image errors above are at most an exponentially small quantity times the square of the particle count. The retained count is at most $T'$, so they remain exponentially small after reducing the exponent $c$. The uniform integrated absolute bound now proves that $F_m(\theta,r)$ is Cauchy at exponential speed in $T'$ along each residue class. Equation~\eqref{eq:normalization-circle-mean} gives \eqref{eq:classwise-normalization} with limits not yet known to be nonzero, since $T'\asymp\log m$.

\paragraph{Saturation excludes a zero averaged limit.}
The individual saturation estimates \eqref{eq:gas-individual-saturation} have exponents $3/8$ at $r=0$ and $5/24$ at $r=i\pi/2$. Thus $\log A_m(r)=o(\log m)$ at both values. If the averaged limit on a residue class were zero, the convergence rate would give $A_m(r)=O(m^{-\delta})$ on that class, a contradiction. Both center values are positive, so every nonzero class limit is positive. There are finitely many classes; taking their minimum and maximum gives bounded factors at both values. Finally $\sigma(0)-\sigma(i\pi/2)=1/6$, proving \eqref{eq:uniform-cylinder-normalization}.
\end{proof}

We next transfer this bounded normalization to plane domains with explicit radius and lattice conventions.

\subsection{Hexagon radii, offsets, and the chord estimate}\label{sec:hexagon-completion}

First consider the Euclidean disk partition $\mathcal W_y(t)$: every polygon surrounds the triangular lattice vertex $y$ and is contained in the open disk $B(y,t)$, with factor $2$ per polygon and the usual critical activity. This partition includes the empty system. For a fixed sufficiently small $c>0$, the uniform geometric comparison gives
\begin{equation}\label{eq:disk-cylinder-comparison}
 \mathcal Z_{m,m}(1)\asymp\mathcal W_y(cN)^2,
 \qquad N=2m.
\end{equation}
Here the cylinder has horizontal period $N$ and its marks are separated by $N/2$. To spell out the comparison, two disk systems of radius $cN$ about the two marks are disjoint and project injectively, giving the lower bound for the cylinder partition. For the upper bound, discard essential polygons and contractible polygons whose specified lifts have diameter at least $cN$. The individual-weight bounds of Section~\ref{c6:uniform-geometry} show that this costs at most a fixed factor: forgetting all avoidance involving a discarded polygon bounds the cost by $\exp(2\sum_L\kappa^{|L|})$. Every remaining polygon lies in the disk of radius $cN$ about the mark it encloses. Indeed a point inside a simple polygon belongs to its convex hull, so its distance from any point of the polygon is at most the polygon's diameter. Separating the two marked types and dropping any avoidance between them gives the upper bound in \eqref{eq:disk-cylinder-comparison}. Changing $c$ by a fixed factor changes only the constants, by the annular bound \eqref{eq:legacy-P-G2}.

Combining \eqref{eq:disk-cylinder-comparison} with the analytic estimate and interpolating monotonically between consecutive sufficiently large even periods gives
\begin{equation}\label{eq:disk-nesting-uniform}
 \mathcal W_y(t)\asymp(1+t)^{1/12}\qquad(t>0),
\end{equation}
uniformly in $y$. For bounded $t$ the lower bound follows from the empty system and the upper bound from a single fixed containing disk. Translation by a triangular lattice vector preserves all weights.

For clarity, the exact pure-sided regular geometry has no unspecified offset. In coordinates
\[
 n_0=(0,1),\quad n_1=(\sqrt3/2,1/2),\quad
 n_2=(\sqrt3/2,-1/2),\quad n_{j+3}=-n_j,
\]
pure lines have levels $d_0\mathbb Z$. If a regular hexagon has center $v$ and inradius $r$, its three support levels $n_j\cdot v+r$ for $j=0,1,2$ must belong to $d_0\mathbb Z$. Since $n_1=n_0+n_2$, the first and third levels minus the second give $r\in d_0\mathbb Z$. Then $n_0\cdot v,n_2\cdot v\in d_0\mathbb Z$, which says exactly that $v$ is a vertex of the triangular lattice. Conversely every $H(y;r)$ in the introduction has all six sides on pure cuts. Its inradius is $r$ and its circumradius is $2r/\sqrt3$.

The disk estimate gives a comparison that also covers nonregular roundings. If $0<a\le b$ are fixed and $B(y,aR)\subset D\subset B(y,bR)$, then positivity gives
\[
 \mathcal W_y(aR)\le Z_D(y)\le\mathcal W_y(bR).
\]
Thus $Z_D(y)\asymp_{a,b}R^{1/12}$ for $R\ge1$. In particular it applies to every pure-support domain
\begin{equation}\label{eq:pure-support-domain}
 D(y;\mathbf b)=\{z:n_j\cdot(z-y)\le b_j,\ 0\le j\le5\},
 \quad b_j\in d_0\mathbb Z,
 \quad aR\le b_j\le bR,
\end{equation}
since it contains $B(y,aR)$ and is contained in $B(y,2bR)$. Some support inequalities can be redundant; the comparison does not require six active sides. Constants depend only on $a,b$, not on the six support levels or the lattice translation. The marked point must remain a fixed fraction of the scale from the boundary for this two-sided bulk estimate.

Here is also the precise effect of translating and rounding a template. Suppose its center is $v$ and its nominal support distances are $r_j$. Replace each absolute level $n_j\cdot v+r_j$ by any $A_j\in d_0\mathbb Z$ with error at most $K d_0$, where $K$ is fixed. Relative to a marked lattice vertex $y$, the resulting levels satisfy
\[
 \bigl|b_j-\bigl(r_j+n_j\cdot(v-y)\bigr)\bigr|\le Kd_0.
\]
Whenever these levels satisfy the fixed inner and outer bounds in \eqref{eq:pure-support-domain}, the same comparison applies, uniformly over all the rounding choices. Thus a size parameter $s$ can replace the actual radius $r$ when $a\le r/s\le b$ for fixed positive $a,b$; equality of the two scales is not needed. The chord threshold below, however, uses the actual inradius.


\begin{proof}[Completion of Theorem~\ref{thm:nesting}]
The nesting assertions follow from \eqref{eq:disk-nesting-uniform} and the domain comparison. We give the chord argument for $H(y;R)$, with $R$ its actual inradius. Proposition~\ref{prop:crossings}, summed over the $O(R)$ boundary ports, bounds the mass of chords with diameter at least $R/2$ by $CR^{3/4}$.

For a lower bound, translate $y$ to zero and use the horizontal caps at heights $-R,R$. Start at the order $R$ ports on the lower cap with horizontal coordinate of absolute value at most $R/10$. Lemma~\ref{lem:narrow-channel-bridges} gives mass at least $cR^{-1/4}$ from each such port for bridges with horizontal displacement at most $R/10$. Their whole trace has $|x|\le R/5$ and $|z_2|\le R$. Consequently its projections on the other four support normals are at most
\[
 (1/2+\sqrt3/10)R<R,
\]
so these bridges lie in the hexagon. Their endpoints are separated vertically by $2R$, hence they are macroscopic. This proves the matching mass lower bound for all sufficiently large admissible $R$.

Every internal link has occupancy mass at most $CR^{1/12}$: the entire hexagon lies in a disk of radius $CR$ around either endpoint of its crossed tiling edge, and Proposition~\ref{prop:occupancy} applies. For a link whose midpoint is in $B(y,R/4)$, both endpoints of that tiling edge lie within $R/4+1/2$ of $y$. For large $R$, disks of radius $R/2$ around them are contained in the hexagon. The same proposition and \eqref{eq:disk-nesting-uniform} therefore give occupancy at least $cR^{1/12}$ at each of the order $R^2$ such links. Every chord using one of them is macroscopic: that midpoint has distance at least $3R/4$ from every boundary point, including its terminal ports.

Sum occupancies over internal links. The upper bound uses all $O(R^2)$ links; the lower bound uses just the central links, all of whose occupying chords meet the diameter restriction. Since a chord with $\ell$ centers uses exactly $\ell-1$ internal links, the first-length mass of the macroscopic chords is comparable to $R^{25/12}$; adding the chord mass accounts for the one-center difference. Dividing by their mass $\asymp R^{3/4}$ proves the claimed mean.

The bounded admissible radii require no asymptotic exception. There are only finitely many such hexagons up to lattice symmetry. Their honeycomb graphs of triangle centers are connected: a segment through the interior between two triangle centers, perturbed to miss tiling vertices, crosses adjacent triangles within the convex hexagon. In particular a shortest dual path from a triangle incident to the lower cap to one incident to the upper cap, with the terminal half-edges added, is a simple chord of diameter at least $2R$. Each small hexagon thus has finite positive macroscopic mass and finite positive first-length mass. Enlarging the constants over these finitely many cases proves the stated bounds for every admissible $R\ge1$.
\end{proof}

The same quantitative argument explains the offset latitude needed later. If the six pure supports satisfy $|b_j-R|\le R/20$, the horizontal cap height is between $19R/10$ and $21R/10$. The preceding tubes satisfy
\[
 n_j\cdot(z-y)\le(21/40+\sqrt3/10)R<19R/20
\]
for the four tilted normals, so the lower chord mass remains uniform. Links with midpoints in $B(y,R/4)$ are at distance at least $7R/10$ from the boundary, while both adjacent lattice vertices have centered disks of radius $R/2$ inside for large $R$. The occupancy and first-length arguments are unchanged. These estimates cover the independent small relative support displacements and bounded row roundings in the sewing construction. They also hold with diameter threshold one half of the domain's actual inradius, which is at most $(b_0+b_3)/2\le21R/20$. For bounded $R\ge1$ there are only finitely many admissible support vectors after translating $y$ to zero. The domain is convex and tiled, and both horizontal caps contain boundary ports. The same connected-dual-graph argument gives an opposite-cap chord of diameter at least $b_0+b_3\ge19R/10$, so its mass and first-length mass are positive for either threshold. Taking the minimum over these finitely many domains extends the estimates to all such bounded scales as well.


\begin{corollary}[Uniform finite chord length]\label{cor:convex-first-length}
Let $D$ be any finite convex pure-sided tiled domain of diameter at most $R\ge1$. Sum over all paths in $D$ with ordered free boundary ports. Then
\[
 \sum_{\gamma\subset D}\ell(\gamma)w(\gamma)\le CR^{25/12},
\]
with an absolute constant, independently of the shape of $D$ and the positions of its ports.
\end{corollary}
\begin{proof}
For every lattice point $y$ incident to $D$, a regular lattice hexagon centered at $y$ of size $CR$ contains $D$. Positivity of polygon weights and Theorem~\ref{thm:nesting} imply $Z_D(y)\le CR^{1/12}$, uniformly even when $y$ is close to the boundary. Proposition~\ref{prop:occupancy}, summed over the $O(R^2)$ internal links, bounds the mass weighted by $\ell-1$ by $CR^{25/12}$. The mass of all chords is at most $CR$: there are $O(R)$ boundary ports, and the positive contour identity bounds the total exit mass from each port by an absolute constant. Adding this mass accounts for the extra center in $\ell$.
\end{proof}
